\subsubsection{A finite double-building sufficient condition}

The decomposition complex of Miller--Patzt--Putman gives a further
conditional form of the two-rank test. This uses homology at two
specific ranks instead of a connectivity conjecture at all ranks.

\begin{proposition}[K17-K12-DOUBLE]\label{prop:k12-double-building}
Let $R=\mathbb Z[1/13!]$, and let $\mathcal T^2(\mathbb Z^N)$
be the double Tits building of compatible flags of proper direct
summands. If
\[
 \widetilde H_{10}(\mathcal T^2(\mathbb Z^9);R)=
 \widetilde H_{11}(\mathcal T^2(\mathbb Z^9);R)=
 \widetilde H_{11}(\mathcal T^2(\mathbb Z^{10});R)=0,
                                                               \tag{DB}
\]
then
\[
 H_4(\mathrm{GL}_9;\mathrm{St}_9\otimes R)=
 H_3(\mathrm{GL}_{10};\mathrm{St}_{10}\otimes R)=0.
\]
Together with the explicit additional hypothesis
$H_5(\mathrm{GL}_8;\mathrm{St}_8\otimes R)=0$, these imply
$K_{12}(\mathbb Z)=0$ by Proposition~\ref{prop:k12-three-ranks}.
The three hypotheses in \textup{(DB)} are not proved here.
\end{proposition}

\begin{proof}
For $i\ge-1$, use the augmented decomposition complex
\[
 S_i(\mathbb Z^N)=
 \bigoplus_{V_1\oplus\cdots\oplus V_{i+2}=\mathbb Z^N}
       \mathrm{St}(V_1)\otimes\cdots\otimes\mathrm{St}(V_{i+2}),
\]
where each direct summand has positive rank. Thus $S_{-1}=\mathrm{St}_N$.
Its differential multiplies adjacent Steinberg apartment classes with
alternating signs. By \cite[Theorem~23.3, p.~49]{MillerPatztPutman2025},
\[
 H_i(S_\bullet(\mathbb Z^N))
   \cong\widetilde H_{i+N-1}(\mathcal T^2(\mathbb Z^N)).
\]
Flat localization preserves this identity.
Theorem~D$'$ of the same source, p.~64, supplies exactness at
$S_{-1},S_0,S_1$ for $N\ge4$. Hence \textup{(DB)} gives exactness
of
\[
 S_d\otimes R\longrightarrow S_{d-1}\otimes R\longrightarrow\cdots
 \longrightarrow S_0\otimes R\longrightarrow\mathrm{St}_N\otimes R
 \longrightarrow0
\]
with $(N,d)=(9,4)$ or $(10,3)$.
The first rank-$9$ hypothesis alone also gives the same exactness
for $(N,d)=(9,3)$; this will recover a needed rank-$9$ boundary
without relying on the announced DEKM endpoint.

Add the kernel of the leftmost differential as a term in degree
$d+1$. This is an exact resolution by $R$-free modules; projectivity
over the group ring is not needed. Its hyperhomology spectral sequence
converges to $H_*(\mathrm{GL}_N(\mathbb Z);\mathrm{St}_N\otimes R)$.
For $0\le p\le d$, its first page is
\[
 E^1_{p,q}=
 \bigoplus_{\substack{r_1+\cdots+r_{p+2}=N\\r_j\ge1}}
 H_q\left(\prod_j\mathrm{GL}_{r_j}(\mathbb Z);
                   \bigotimes_j(\mathrm{St}_{r_j}\otimes R)\right).
                                                               \tag{DB1}
\]
Indeed, $\mathrm{GL}_N(\mathbb Z)$ is transitive on ordered integral
decompositions with fixed ranks, their stabilizers are the displayed
products, and Shapiro's lemma applies. This is the single-rank
construction in the complete proof of
\cite[Lemma~23.5, pp.~50--51]{MillerPatztPutman2025}.
The extra kernel term is in column $d+1$ and cannot occur in total
degree $d$.

The low-rank homology data in
\cite[Theorems~4.3 and~8.1]{ElbazGanglSoule2013}, together with
the documented cone ranks through $12$ at rank $8$ in
\cite[Proposition~1.5, Tables~3--5]{DEKM2026}, localized in $R$, say
that for $1\le r\le8$ and $0\le q\le4$ the only nonzero groups
are
\[
 H_0(\mathrm{GL}_1;\mathrm{St}_1\otimes R)=R,
 \qquad
 H_3(\mathrm{GL}_3;\mathrm{St}_3\otimes R)=
 H_3(\mathrm{GL}_4;\mathrm{St}_4\otimes R)=R.
\]
Here the rank-one statement also follows directly from
$\mathrm{GL}_1(\mathbb Z)=C_2$. These ranges require only cone
rank $12$ at rank $8$, hence are within the displayed enumeration.
The additional rank-$9$ degree-three endpoint would require cone
rank $12$ there, beyond the posted rank-$11$ enumeration; it is
derived next from \textup{(DB)} instead. Every module in the
displayed low-rank ranges is free over $R$, so no Tor terms occur
in the required iterated K\"unneth formulas.

First take $(N,d)=(9,3)$. The hypothesis
$\widetilde H_{10}(\mathcal T^2(\mathbb Z^9);R)=0$ gives
exactness at $S_2$. Thus the partial resolution starts at $S_3$,
and its additional kernel lies in degree $4$. Every factor in
\textup{(DB1)} has rank at most $8$. On $p+q=3$, a nonzero
K\"unneth term could have only rank-one factors, giving
$q=0,p=3$ and total rank $5$, or one rank-three or rank-four
factor in degree three, giving $q=3,p=0$ and total rank at most
$4+1=5$. Neither can give rank $9$. Therefore
\[
 H_3(\mathrm{GL}_9;\mathrm{St}_9\otimes R)=0. \tag{DB2}
\]
The lower rank-$9$ groups vanish by the documented cone-rank
range through $11$ (and independently the degree-one and
degree-two groups vanish by Corollary~E of
\cite{MillerPatztPutman2025}).

For $(N,d)=(9,4)$, a nonzero K\"unneth contribution to
\textup{(DB1)} on $p+q=4$ could have only rank-one factors,
or one rank-three or rank-four factor in degree three and rank-one
factors otherwise. In the first case $q=0,p=4$ and the total rank
is six. In the second $q=3,p=1$ and the total rank is at most
$4+1+1=6$. Neither has rank nine.

For $(N,d)=(10,3)$ the rank-nine factors contribute nothing through
degree three, now using \textup{(DB2)}. The same list of possibilities applies. Rank-one
factors alone give $q=0,p=3$ and total rank five. A single
rank-three or rank-four factor gives $q=3,p=0$ and total rank at
most $4+1=5$. Neither has rank ten.

Thus every first-page term in the relevant total degree is zero in
both spectral sequences. They give the required degree-four and
degree-three Steinberg vanishings. Under the additional explicit
rank-eight hypothesis, Proposition~\ref{prop:k12-three-ranks}
gives the final conditional conclusion. The announced rank-eight
endpoint is not independently inferred from the posted enumeration.
\end{proof}
